\clearpage
\fancyhead[L]{\small BODE. ASTRONOMISCHES JAHRBUCH.}
\hypertarget{native-635-lower}{}
\begin{center}\small Göttingische gelehrte Anzeigen. Stück 2. Seite 17..24. 1824 Januar 3.\\[3pt]\rule{.72\linewidth}{.3pt}\end{center}

\emph{Astronomisches Jahrbuch für das Jahr \textup{1826} nebst einer Sammlung der neuesten in die astronomischen Wissenschaften einschlagenden Abhandlungen, Beobachtungen und Nachrichten. Berechnet und herausgegeben von} Dr. J. E. \textsc{Bode}, \emph{Königl. Astronom u. s. w.} Berlin 1823. Bei dem Verfasser und in Commission bei Ferd. Dümmler. 256 Seiten in Octav, nebst einer Kupfertafel.

Mit diesem Jahrgange fängt das zweite halbe Hundert in einer Reihe an, wodurch zur Verbreitung astronomischer Kenntnisse in Deutschland viel beygetragen ist, und der wir noch langen Fortgang \SourceLeaf{636} unter der Leitung des würdigen Herausgebers wünschen. Unter den Ereignissen am Himmel im Jahre 1826, die ein allgemeineres Interesse haben, zeichnen wir eine totale Mondfinsterniss (am 14. November), eine partiale Sonnenfinsterniss (am 29. November) und eine Bedeckung des Saturn vom Monde (am 16. Februar) aus.

Die Abhandlungen und Nachrichten im Anhange nehmen 168 Seiten ein. Zuerst einige Resultate aus der Triangulirung im Königreiche Hannover von Hrn. Hofr. \textsc{Gauss}. Es wird hier zuerst die geographische Lage von 22 ausgewählten Punkten mitgetheilt, wie sie, aus dem vorläufigen Anschlusse an die von \textsc{Zach}'sche Basis, im Herbst 1822 abgeleitet ist. Wir können hier die Bemerkung hinzufügen, dass der seitdem ausgeführte Anschluss an die von Hrn. \textsc{Schumacher} in Holstein gemessene Grundlinie, diese Resultate gar nicht bemerkbar abändert; nur die nordlichern Breiten werden ein Zehntheil einer Secunde vermindert. Die drei Thürme in Hildesheim, welche in diesem Verzeichniss ohne Namen vorkommen, können jetzt namhaft gemacht werden: der erste ist der Andreas-, der zweite der Michaelis-, der dritte der Jacobi-Thurm. Dann die Uebersicht der aus den gegenseitigen Zenithdistanzen auf 28 Linien in dem Dreieckssystem gefolgerten terrestrischen Refractionen. Das Mittelverhältniss der Krümmung des terrestrischen Bogens zu der beobachteten ganzen Refraction ist wie 1 zu 0.1306. Durch die im Jahr 1823 neuhinzugekommenen Messungen wird das Resultat noch ein wenig kleiner, nemlich wie 1 zu 0.1278. Dieses Resultat ist bedeutend kleiner, als das Verhältniss, welches man bisher anzunehmen pflegte, wo man die halbe Refraction (gewöhnlich, obgleich nicht ganz richtig, schlechtweg Refraction genannt) zu 0.08 der Krümmung des irdischen Bogens ansetzte. Inzwischen scheint jene Bestimmung, wo bei den Beobachtungen fast durchgängig Heliotroplicht den Zielpunkt abgegeben hat, und wo so viele grosse Distanzen vorkommen (die grösste ist $11\tfrac{1}{2}$ Meile) mehr Zutrauen zu verdienen, als andere, bei denen das Visiren auf terrestrische Zielpunkte schon an sich keine so grosse Genauigkeit in den beobachteten Zenithdistanzen verstattete, und wo überdies die Entfernungen gewöhnlich viel kleiner waren. Bei kleinen Entfernungen sind, auch nach diesen Erfahrungen, die Anomalien immer am grössten, und besonders in den Vormittagsstunden, um Mittag und in den frühern Nachmittagsstunden, wo, jenen zufolge, an sonnigen Sommertagen in flachen Gegenden und bei kleinern Entfernungen, das nahe über den Erdboden wegstreichende Licht gewöhnlich eine negative Refraction erleidet. --- Beiträge zu geographischen Längenbestimmungen aus beobachteten Sonnenfinsternissen und Sternbedeckungen von Hrn. Prof. \textsc{Wurm}. --- Beobachtungen von Planeten-Oppositionen, Verfinsterungen der Jupiterstrabanten und Sternbedeckungen, angestellt auf der Sternwarte in Wilna von Hrn. \textsc{Sniadecky}. --- Verschiedene Beobachtungen des Hrn. \textsc{Rümker} in Paramatta: die wichtigsten derselben, nemlich die des \textsc{Enke}'schen Cometen im Juni 1822, sind bereits im 26. Stück dieser Blätter Jahrgang 1823 [S. 442 d. B.] mitgetheilt. --- Der hierauf folgende interessante Aufsatz des Hrn. Dr. \textsc{Olbers} berührt die erhabenste Seite der Astronomie, die Frage nemlich, ob das Weltall Grenzen habe, oder nicht. Bei dem Bewusstsein unsers Unvermögens, den Schleier aufzuheben, welcher dieses Geheimniss deckt, beschäftigen wir uns doch gern mit der Aufklärung solcher Zweifel, welche die Annahme, das Weltall sei unendlich, mit unsern sonstigen Einsichten in Widerspruch zu stellen scheinen könnten. Auf einen solchen Zweifel, wenigstens an der unendlichen Anzahl von Sonnen, führt die Erwägung, dass, bei Annahme derselben, \emph{jede} vom Auge ausgehende Richtung irgendwo in kleinerer oder grösserer Ferne auf einen Sonnenkörper treffen, und daher das ganze Himmelsgewölbe in nirgends unterbrochenem Sonnenglanz erscheinen würde. Dieser Schluss ist an sich vollkommen richtig: allein ein gleichförmiger Glanz würde nur dann Statt finden können, wenn der Weltraum absolut durchsichtig wäre, und also das Licht, indem es ihn durchdringt, gar nichts an Intensität verlöre. Diese Voraussetzung ist aber durchaus unerwiesen und in der That schon an sich höchst unwahrscheinlich. Hr. Dr. \textsc{Olbers} zeigt nun, dass eine \SourceLeaf{637} verhältnissmässig nur sehr geringe Undurchsichtigkeit angenommen zu werden braucht, um den scheinbaren Widerspruch der Erfahrung mit einer unendlichen Zahl von Sonnen aufzuheben. Ist die Schwächung des Lichts auf seinem Wege vom Sirius zu uns nur $\frac{1}{800}$, so würde eine hohle Kugel mit einem Halbmesser von, 30000 Siriusweiten, in deren Mitte wir uns befänden und deren ganze innere Fläche Sonnenglanz ausströmte, uns nur $\frac{1}{700000}$ so hell erscheinen, als der Himmelsgrund in einer heitern Vollmondsnacht, also gewiss gar keinen merklichen Eindruck auf unsere Augen hervorbringen. --- Fixsternverzeichniss von Hrn. \textsc{Pond}, mitgetheilt von Hrn. Prof. \textsc{Tralles}. Der Sinn dieses Verzeichnisses und der daraus gezogenen Folgerung ist uns nicht ganz verständlich. --- Fortgesetzte Nachrichten über den \textsc{Pons}'schen (\textsc{Enke}'schen) Cometen von Hrn. Prof. \textsc{Enke}. Die neuen Beobachtungen dieses Weltkörpers im Jahre 1822 bestätigen das schon aus den frühern folgende merkwürdige Phänomen, dass neben den durch die Planetenstörungen erklärbaren Veränderungen der Umlaufszeit noch eine successive Abnahme derselben, die von einem Umlauf zum andern jetzt etwas über drei Stunden beträgt, Statt findet. Wahrscheinlich haben wir hierin die erste sichere Erfahrung von dem wirklichen Dasein eines Widerstandes, welchen die Cometen bei ihrer Bewegung im Weltraume erleiden. Hoffentlich wird dieser Weltkörper schon bei seiner nächsten Annäherung zur Sonne im August 1825 wieder in Europa beobachtet werden können; wenigstens wird über den Platz, wo er zu suchen sein wird, nach der Bearbeitung des eben so geschickten als unermüdeten Rechners, gar keine Ungewissheit Statt finden. --- Astronomische Beobachtungen auf der Prager Sternwarte von den Hrn. \textsc{David} und \textsc{Bittner}; ähnliche Beobachtungen in Kremsmünster von Hrn. \textsc{Derflinger} und in Prag von Hrn. \textsc{Hallaschka}. --- Originalbeobachtungen des dritten Cometen von 1822 von Hrn. Dr. \textsc{Olbers}. --- Einige mechanische Untersuchungen über die Entstehungen der Cometenschweife von Hrn. Dr. \textsc{Lehmann}. Dass alle Versuche, diesen räthselhaften Gegenstand befriedigend zu erklären, bisher ohne erheblichen Erfolg gewesen sind, darf uns nicht wundern. Erfahrungen von einer solchen Bestimmtheit, dass daraus scharfe Resultate zur Grundlage strenger Rechnungen gewonnen werden könnten, fehlen uns gänzlich; über die Beschaffenheit der dabei thätigen Kräfte wissen wir wenig gewisses, und die strenge mathematische Behandlung des Erfolgs von hypothetisch angenommenen Kräften hat sehr grosse Schwierigkeiten. Dieser Versuch des Hrn. \textsc{Lehmann}, die Cometenschweife blos aus den bekannten Kräften, ohne Zuziehung von Repulsivkräften zu erklären, bleibt zwar (wie alle andern versuchten Erklärungen) auch noch viel zu sehr an der Oberfläche des Problems, und hält sich noch viel zu sehr an inadäquate Vorstellungen im Allgemeinen, als dass sich so auch nur über die Möglichkeit einer Erklärung auf diesem Wege absprechen liesse; indessen ist er nicht ohne sinnreiche Ansichten, und man kann nicht läugnen, dass jenes auch von den bisherigen Versuchen mit Repulsivkräften gilt, deren Nothwendigkeit noch nicht als entschieden betrachtet werden darf. --- Astronomische Beobachtungen auf der königl. Sternwarte in Berlin im Jahre 1822. --- Ein stärker vergrössernder Ocular-Ansatz für achromatische Fernröhre, erfunden von Hrn. W. \textsc{Kitchner} in London, So viel man aus der sehr unvollständigen und unklaren Angabe schliessen kann, besteht die Erfindung nur darin, die Oculargläser gegen einander beweglich zu stellen, um damit veränderliche Vergrösserungen zu erhalten. --- Beobachtungen und Elemente des Cometen vom September 1822 von Hrn. \textsc{Rümker} in Paramatta. Diese Beobachtungen gehen bis gegen die Mitte des November, also beträchtlich weiter, als die Europäischen, und verdienen daher, mit diesen verbunden zu werden. Ferner Beobachtung des Merkur-Durchganges am 5. November 1822 von demselben. --- Geographische Ortsbestimmungen in der Altmark und anderen Grenzen von Hrn. Musik-Director \textsc{Stöpel} in Tangermünde. Dieses schätzbare Verzeichniss enthält 134 Ortsbestimmungen nach trigonometrischen Messungen, die sich an die Dreiecksseite in der grossen von \textsc{Müffling}'schen Triangulirung von Magdeburg zum Hagelsberg anschliessen. Wir hätten gewünscht, dass Hr. \textsc{Stöpel} zugleich einiges über seine Hülfsmittel und seine \SourceLeaf{638} Beobachtungsart mitgetheilt und sich nicht blos deswegen auf ein auswärts unbekanntes Provinzialblatt bezogen hätte. --- Beobachtung einer Sternbedeckung von Mars, von Hrn. Pr. \textsc{Tralles}, um so merkwürdiger, je seltener solche Erscheinungen sind. Der Stern wurde nach seinem Austritt erst in einer Entfernung von mehr als einem Planetendurchmesser wieder sichtbar, anfangs sehr schwach, dann aber schnell an Licht zunehmend; das Fernrohr hatte $3\tfrac{1}{2}$ Zoll Oeffnung. Hr. \textsc{Tralles} schliesst daraus auf das Dasein einer bedeutenden Marsatmosphäre. --- Ueber die von Hrn. G. R. \textsc{Pastorf} entdeckte Photosphäre der Planeten von Hrn. \textsc{Ritz}. Hr. \textsc{Ritz} erklärt diese Erscheinung aus einer doppelten Reflexion zwischen den Objectivlinsen. Rec. hat immer dieselbe Ansicht gehabt und sich durch Rechnung nach den wirklichen Dimensionen von \textsc{Frauenhofer}'schen Fernröhren überzeugt, dass diese durch die gedachte Reflexion ein solches Phänomen, wie das beobachtete ist, erzeugen müssen. Eine wesentliche Bedingung der Entstehung ist, dass die Krümmungen der beiden einander zugekehrten Flächen \emph{beinahe} gleich sein müssen, wie es bei der Construction der \textsc{Frauenhofer}'schen Objective wirklich der Fall ist. --- Bemerkungen über den vorigen Gegenstand von Hrn. Justiz-Commissionsrath \textsc{Kunovsky}. Hr. \textsc{Kunovsky} gibt der Erklärungsart des Hrn. \textsc{Ritz}, dessen Aufsatz vor dem Abdruck ihm zur Einsicht mitgetheilt zu sein scheint, seinen Beifall; Rec. gesteht indessen, dass ihm vorgekommen ist, als habe Hr. \textsc{Kunovsky} den Geist der Erklärung nicht richtig aufgefasst, da beinahe alles was er als eine Bestätigung anführt, wenn man es als factisch bewiesen ansehen dürfte, in Widerspruch damit stehen oder wenigstens einer ganz anderen besondern Erklärung bedürfen würde. Hr. \textsc{Kunovsky} behauptet, die Erscheinung mit \emph{allen} Fernröhren, die er auf die Probe gestellt hat, bemerkt zu haben, obgleich bekannt ist, dass die englischen Objective ganz anders construirt sind, als die \textsc{Frauenhofer}'schen; er findet die Lichtsphäre bei verschiedenen Fernröhren nahe der Oeffnung des Objectivs proportional, was im Geist der obigen Erklärung nur von \emph{Einem und demselben} Fernrohre gilt, insofern dem Objectiv bald eine grössere, bald eine geringere Oeffnung gelassen wird; bei verschiedenen Fernröhren, deren Dimensionen ungleich, aber in \emph{allen} Stücken einander proportional wären, müsste die Lichtsphäre durchaus gleich gross erscheinen; endlich hat Hr. \textsc{Kunovsky} bei Anwendung von viel stärkern Vergrösserungen die Lichtsphäre des Jupiter, wie es sein muss, viel blässer werden sehen, während die Lichtsphäre von $\alpha$ in der Leyer durch stärkere Vergrösserung fast gar nichts von ihrer Helligkeit verloren haben soll, was auch wiederum mit obiger Erklärung unverträglich sein würde. Sehr brauchbar ist übrigens die von Hrn. \textsc{Kunovsky} angedeutete Prüfung der Reinheit des Glases der Objective. --- Aus einem Schreiben des Hrn. Prof. \textsc{Bürg} werden einige Resultate seiner neuern Untersuchungen über die Mondstheorie mitgetheilt. --- Fragmente zur Erklärung des Aratus, von Hrn. Prof. \textsc{Schaubach}. --- Beobachtung der totalen Mondfinsterniss am 26. Januar 1823 und einer Sternbedeckung von Hrn. \textsc{Rümker} im Paramatta. --- Einige Beobachtungen auf der Dorpater Sternwarte. --- Aehnliche auf der Sternwarte in Krakau von Hrn. Prof. \textsc{Leski}. --- Ueber die astronomische Strahlenbrechung von Hrn. Prof. \textsc{Bessel}. --- Astronomische Nachrichten von Hrn. Prediger \textsc{Luthmer} in Hannover; \textsc{Herschel}'s Grabschrift. --- Neue Elemente der Junobahn von Hrn. Prof. \textsc{Nicolai}. Dies sind die sehr schätzbaren Resultate einer mühsamen, auf Quadraturen gegründeten Berechnung der Störungen, welche Jupiter seit 1804 auf die Bewegung der Juno ausgeübt hat, und der Discussion von 15 seit der Entdeckung beobachteten Oppositionen. Diese werden durch die Rechnung mit vieler Genauigkeit dargestellt. Die noch übrig bleibenden kleinen Unterschiede ist Hr. \textsc{Nicolai} geneigt einer von ihm vermutheten Unzulänglichkeit des \textsc{Newton}'schen Gravitationsgesetzes zuzuschreiben, indem er es nemlich für möglich hält, dass bei gleicher Entfernung Jupiter eine andere Attraction auf die Sonne als auf die Juno ausübt. Für absolut unmöglich kann man freilich eine solche auch schon von andern aufgestellte Hypothese nicht erklären, allein sehr unwahrscheinlich ist sie doch schon deswegen, weil es factisch bewiesen ist, dass die Attraction der Sonne auf alle Haupt- und Nebenplaneten \SourceLeaf{639} gleich gross ist, eben so wie die Gravitation der verschiedensten Körper auf der Erdoberfläche gegen die Erde: nach Rec. Ansicht, darf man daher zu einer der Analogie so sehr widersprechenden Voraussetzung nur dann erst seine Zuflucht nehmen, wenn dafür entscheidende Beweise vorhanden sind, was bis jetzt keinesweges der Fall ist. --- Beobachtete Sternbedeckungen auf der Wiener Sternwarte von Hrn. Prof. \textsc{Littrow}. --- Verzeichniss von 795 Doppelsternen aus Hrn. \textsc{Struve}'s Astr. Beobachtungen, Band 3. --- Unter den kürzeren astronomischen Nachrichten und Bemerkungen findet sich auch noch manches interessante, obwohl sich über verschiedene Artikel noch Erinnerungen machen liessen. So sieht z. B. Rec. nicht recht ein, in wie fern der Umstand, dass bei dem \textsc{Halley}'schen Cometen die Unterschiede der beobachteten Umlaufszeiten \emph{untereinander} viel grösser sind, als bei dem \textsc{Enke}'schen Cometen, einen Beweis von den grossen Fortschritten der Beobachtungskunst und des Calculs geben soll, so unläugbar letztere an sich sind. Selbst dass der Unterschied zwischen der \emph{berechneten} Umlaufszeit und der beobachteten bei dem \textsc{Halley}'schen Cometen viel grösser war, als bei dem \textsc{Enke}'schen, beweist an sich noch wenig für die Fortschritte des Calculs, da die Umstände bei beiden Cometen so sehr verschieden waren.
